\section{Dilution: a correct argument about normalisation, and the opposite measurement}
\label{sec:dilution}

The most-repeated qualitative argument for state space models in vision
concerns large objects. Self-attention weights sum to one, so an object
spanning $K$ tokens receives an average weight of $1/K$; the more of the
image the object occupies, the more thinly the attention mass is spread over
it. A state space model has no such constraint---each object token adds to
the state as it is scanned---and is therefore said not to suffer the
dilution.

The argument about normalisation is correct. The conclusion drawn from it is
not the one the data supports.

\begin{figure}[tbp]
  \centering
  \includegraphics[width=\linewidth]{fig3_dilution.pdf}
  \caption{Object coverage by area, for pretrained and untrained weights.
  What is plotted is the excess over the random baseline; raw precision@K
  rises with object area for every model and is not interpretable on its
  own.}
  \label{fig:dilution}
\end{figure}

\begin{table}[tbp]
\centering
\caption{Above-baseline precision@K by object area, pretrained weights. Vim-S is the lowest in every bin. Raw precision@K rises with object area for all three models because the random baseline rises faster; only the excess over that baseline is interpretable.}
\label{tab:e3}
\begin{tabular}{lrrr}
\toprule
Area bin & DeiT-S & CMT-S & Vim-S \\
\midrule
<2\% & 0.4054 & 0.2660 & 0.1682 \\
2-5\% & 0.3574 & 0.2636 & 0.1602 \\
5-10\% & 0.3325 & 0.2669 & 0.1920 \\
10-20\% & 0.2566 & 0.1888 & 0.1398 \\
20-40\% & 0.2239 & 0.1627 & 0.1130 \\
>=40\% & 0.0782 & 0.0569 & 0.0366 \\
\bottomrule
\end{tabular}
\end{table}


\subsection{How this must be read, and how it is easy to misread}

Raw precision@K \emph{rises} with object area for all three models. That is
not evidence against dilution; it is arithmetic. As an object grows, so does
$K$, and a random ranking would also place more of its mass inside the
object. The random baseline rises faster than the models do, and the quantity
that carries information is the excess over that baseline.

A reader who quotes raw precision@K without the baseline will describe this
experiment as saying the opposite of what it says. Table~\ref{tab:e3}
therefore reports only the excess.

One further caution: the six area bins hold different object instances drawn
by stratified sampling, so the decline across bins is not a paired
comparison. Differences between bins should not be read as a within-object
trajectory.

\subsection{The measurement}

Vim-S has the lowest excess over baseline in every one of the six area bins.
Not at large object sizes only, where the dilution argument makes its
prediction---everywhere.

This is what makes the result scale-independent, and the distinction matters
because our registered prediction concerned the \emph{slope}. Measured as a
ratio, the excess of all three models falls by a similar factor from the
smallest bin to the largest, which would say the models degrade alike.
Measured as an absolute drop, Vim-S loses the least of the three, which is
exactly the shape the prediction described---but only because it starts from
the lowest value and so has the least to lose. The choice of scale produces
the verdict, so we do not report one. We place the finding at the level
instead, where no such choice is available: at every object size measured,
the state space model puts less attribution on the object than either
attention model.

\subsection{Why the argument does not reach the conclusion}

The dilution argument establishes that a state space model is free of a
particular normalisation constraint. It does not establish that the model
places more of its attribution on the object, and those are different claims.

Two mechanisms stand in the way, and the earlier comparison we are auditing
names the first of them itself, in a later section, without connecting it to
this one: the recurrent state is a narrow channel of a few tens of
dimensions. Absence of a normalisation constraint does not create capacity.
If the state cannot hold the object's information, it does not matter that
nothing is dividing that information by $K$.

The second is the scan order. Under a row-major flattening the tokens of a
two-dimensional object are not contiguous in the sequence; they are
interleaved with background between row crossings. What the state
accumulates as it passes over an object is therefore repeatedly interrupted
by material the object does not contain.

\subsection{An independent confirmation from a different measurement}

If the scan order is implicated, then objects that are tall should fare worse
than objects that are wide, since a tall object crosses more rows and is
therefore more fragmented in the sequence. That is a prediction about a
quantity we measured for another reason.

It holds, and only for the state space model. The difference between tall and
wide objects in precision@K is \DeitTallWide{} for DeiT-S ($z =
\DeitTallWideZ{}$) and \CmtTallWide{} for CMT-S ($z = \CmtTallWideZ{}$),
both indistinguishable from zero. For Vim-S it is \VimTallWide{} ($z =
\VimTallWideZ{}$). Without training the gap widens to
\VimTallWideRandom{} ($z = \VimTallWideRandomZ{}$).

The direction, the isolation to one model, and the widening in the untrained
condition all match what Section~\ref{sec:erf} measured about receptive field
shape. Two experiments built for different purposes, using different data and
different metrics, point at the same structural cause.
